% ============================================================
%  ch13 —— §9（原文 p.48–p.52）
% ============================================================
\routenote{p.48--p.52}

\section{附录：几何 Satake 等价与 Langlands 对偶群的钉住}

先需要张量函子的若干一般性质。设 $\mathcal C$ 是带纤维函子 $\omega:\mathcal C\to\Vec_F$（$F$ 是域）的刚性张量范畴，$H=\operatorname{Aut}^\otimes(\omega)$ 为由 $(\mathcal C,\omega)$ 决定的 $F$ 上代数群。

对任何 $F$-代数 $R$，令 $\omega_R:\mathcal C\xrightarrow{\omega}\Vec_F\xrightarrow{\otimes_FR}\mathrm{Mod}_R$。记 $\operatorname{Aut}^\otimes(\mathcal C,\omega_R)$ 为对 $(\sigma,\alpha)$ 的同构类构成的集合，其中 $\sigma:\mathcal C\xrightarrow{\sim}\mathcal C$ 是张量自等价，$\alpha:\omega_R\circ\sigma\Rightarrow\omega_R$ 是函子的自然同构。则 $\operatorname{Aut}^\otimes(\mathcal C,\omega_R)$ 带自然的群结构。记 $\operatorname{Aut}^\otimes(\mathcal C,\omega)$ 为函子 $R\mapsto\operatorname{Aut}^\otimes(\mathcal C,\omega_R)$，它定义 fppf 群层。另一方面，记 $\operatorname{Aut}(H)$ 为 $F$ 上 pro-代数群 $H$ 的自同构 fppf 层。

\begin{lemma}
存在 fppf 群层的自然同构 $\operatorname{Aut}^\otimes(\mathcal C,\omega)\xrightarrow{\sim}\operatorname{Aut}(H)$。特别地，有群的同构 $\operatorname{Aut}^\otimes(\mathcal C,\omega)\xrightarrow{\sim}\operatorname{Aut}(H)$，它诱导群的同构 $[\operatorname{Aut}^\otimes(\mathcal C)]\xrightarrow{\sim}\operatorname{Out}(H)$，其中 $[\operatorname{Aut}^\otimes(\mathcal C)]$ 是 $\mathcal C$ 的张量自等价的同构类集合。
\end{lemma}

在 $F$-点的层面，对 $(\sigma,\alpha)\in\operatorname{Aut}^\otimes(\mathcal C,\omega)$，它给出 $H=\operatorname{Aut}^\otimes(\omega)$ 的下述自同构：把 $h:\omega\Rightarrow\omega$ 送到自然变换
\[
\omega\xrightarrow{\alpha}\omega\circ\sigma\xrightarrow{h\circ\mathrm{id}_\sigma=\mathrm{id}}\omega\circ\sigma\xrightarrow{\alpha^{-1}}\omega.
\]

更一般地，设给定张量函子 $\Phi:\mathcal C\to\mathcal C'$，其中 $\mathcal C'$ 是另一刚性张量范畴，则可类似地定义群层 $\operatorname{Aut}^\otimes(\mathcal C,\Phi)$。

设 $\omega':\mathcal C'\to\Vec$ 是纤维函子，$\omega=\omega'\circ\Phi$。则把 $(\sigma,\alpha)\in\operatorname{Aut}^\otimes(\mathcal C,\Phi)$ 送到 $(\sigma,\mathrm{id}_{\omega'}\circ\alpha)\in\operatorname{Aut}^\otimes(\mathcal C,\omega)$，给出自然同态 $\omega'_*:\operatorname{Aut}^\otimes(\mathcal C,\Phi)\to\operatorname{Aut}^\otimes(\mathcal C,\omega)=\operatorname{Aut}(H)$。换言之，$\operatorname{Aut}^\otimes(\mathcal C,\Phi)$ 作用在 pro-代数群 $H$ 上。另一方面，有 pro-代数群的自然同态 $\Phi^*:H'=\operatorname{Aut}^\otimes(\omega')\to H=\operatorname{Aut}^\otimes(\omega)$。

\begin{lemma}
同态 $\Phi^*:H'\to H$ 经 $H'\to H^{\operatorname{Aut}^\otimes(\mathcal C,\Phi)}\subset H$ 分解。
\end{lemma}

现考虑 §2.3 中标准化的半单 Satake 范畴 $\mathcal S$。沿袭 [30] 与 [18]，用整体截面函子 $h$ 定义 $\ckG=\operatorname{Aut}^\otimes(h)$，并得到几何 Satake 等价 $\mathcal S\cong\Rep(\ckG)$。

\begin{lemma}
\begin{enumerate}[label=(\arabic*), leftmargin=1.9em, itemsep=.2em]
\item 存在自然同态 $\operatorname{Aut}(G)\to\operatorname{Aut}^\otimes(\mathcal S,h)\cong\operatorname{Aut}(\ckG)$，它经 $\tilde\iota:\operatorname{Out}(G)\to\operatorname{Aut}(\ckG)$ 分解；
\item 存在 $\ckG$ 的自然钉住 $\dagger=(\check B,\check T,\{x_{\alpha_i^\vee}\})$，它被经 $\tilde\iota$ 的 $\operatorname{Out}(G)$-作用保持。记 $\operatorname{Aut}^\dagger(\ckG)$ 为 $\ckG$ 的保持该钉住的自同构群。则 $\tilde\iota$ 诱导同构 $\iota:\operatorname{Out}(G)\xrightarrow{\sim}\operatorname{Aut}^\dagger(\ckG)$。
\end{enumerate}
\end{lemma}

\begin{proof}
(1) $\operatorname{Aut}(G)$ 在 $(\mathcal S,h)$ 上的作用由其在对 $\mathrm{Gr}_G$ 的作用诱导。由于 $\mathcal S$ 中的对象带 $G$ 在 $\mathrm{Gr}_G$ 上共轭作用的 $G$-等变结构，$G$ 的内自同构在 $(\mathcal S,h)$ 上平凡作用，即 $\operatorname{Aut}(G)$ 在 $(\mathcal S,h)$ 上的作用经 $\operatorname{Out}(G)$ 分解。

(2) 先需给出一个被 $\operatorname{Aut}(G)$-作用保持的 $\ckG$ 的钉住。为此需给出极大环面 $\check T\subset\ckG$、余特征 $2\rho\in X_*(T)=X_*(\check T)$（钉住中正余根和之半）与主幂幺元 $e\in\ckg$（单根向量之和）。等价地，需要
\begin{itemize}[leftmargin=1.6em, itemsep=.2em]
\item[(P1)] 把纤维函子 $h$ 分解为张量函子
\[
h=\bigoplus_{\mu\in X_*(T)}h_\mu:\mathcal S\to\Vec^{X_*(T)}\xrightarrow{\text{遗忘}}\Vec;
\]
\item[(P2)] 找张量导子 $e:h\to h$（即 $e(K_1*K_2)=e(K_1)\otimes\mathrm{id}_{K_2}+\mathrm{id}_{K_1}\otimes e(K_2)$），它把 $h_\mu$ 送到 $\oplus_ih_{\mu+\alpha_i}$，且对每个单根 $\alpha_i$，分量 $h_\mu\to h_{\mu+\alpha_i}$ 作为函子非零。
\end{itemize}

分解 $h_\mu$ 由 Mirković--Vilonen 的“权函子”给出，见 [30, Theorem 3.5, 3.6]。他们还证明 $h=\bigoplus_ih_i$（$h_i$ 是满足 $\langle2\rho,\mu\rangle=i$ 的 $h_\mu$ 之和）与 $h=H^*(\mathrm{Gr}_G,-)$ 的上同调次数分次一致。张量导子 $e$ 由与 $c_1(\mathcal L_{\det})\in H^2(\mathrm{Gr}_G,\mathbb{Q}_\ell)$ 的 cup 积给出，其中 $\mathcal L_{\det}$ 是 $\mathrm{Gr}_G$ 上的行列式线丛。

$\operatorname{Aut}^\dagger(G)\subset\operatorname{Aut}(G)$ 在 $\mathcal S$ 上的作用以与置换 $\mu\in X_*(T)$ 相同的方式置换权函子 $h_\mu$，保持 $\langle2\rho,\mu\rangle$（从而保持 $2\rho\in X_*(T)=X_*(\check T)$），并与 $e=c_1(\mathcal L_{\det})$ 交换。因此 $\operatorname{Aut}^\dagger(G)\xrightarrow{\sim}\operatorname{Out}(G)$ 在 $\mathcal S$ 上的作用保持上述钉住。

$\sigma\in\operatorname{Aut}^\dagger(G)$ 诱导 $\ckG$ 的 Dynkin 图的对偶自同构 $\check\sigma$。由于 $\sigma^*\mathrm{IC}_\mu\cong\mathrm{IC}_{\sigma^{-1}(\mu)}$，$\mathcal S\cong\Rep(\ckG)$ 的自等价 $\sigma^*$ 同构于 $\Rep(\ckG)$ 的由 $\ckG$ 的 Dynkin 图上对偶自同构 $\check\sigma^{-1}$ 给出的钉住自同构所诱导的自等价。这证明
\[
\operatorname{Out}(G)\to[\operatorname{Aut}^\otimes(\mathcal S)]\cong[\operatorname{Aut}^\otimes(\Rep(\ckG))]\cong\operatorname{Out}(\ckG)
\]
是同构（最后一个同构来自引理 9.1）。因此 $\iota:\operatorname{Out}(G)\to\operatorname{Aut}^\dagger(\ckG)$ 也是同构。
\end{proof}

为了 §6 的目的，还需要 (2.11) 中定义的另一纤维函子 $\omega_\varphi$。回忆 $S_\varphi$ 是 $\varphi$ 在 $T\rtimes\operatorname{Aut}^\dagger(G)\times\mathbb{G}_m^{\mathrm{rot}}$ 作用下的稳定子，$S^1_\varphi(k)=\ker(S_\varphi(k)\to\mathbb{G}_m^{\mathrm{rot}}(k))$。由 (6.4) 回忆 $S^1_\varphi(k)$ 作用在 ${}_\varphi\ckG$ 上。

\begin{lemma}
存在 ${}_\varphi\ckG$ 的自然钉住 $\ddagger=({}_\varphi\check B,{}_\varphi\check T,\{{}_\varphi x_{\alpha_i^\vee}\})$，它被 $S^1_\varphi(k)$-作用保持。
\end{lemma}

\begin{proof}
利用 (4.4)，可把纤维函子 $\omega_\varphi$ 改写为
\[
\mathcal S\xrightarrow{\Psi}\mathrm{Perv}_{I_0}(\mathrm{Fl}_G)\xrightarrow{V_\varphi}\Vec,
\]
其中 $\mathrm{Fl}_G=G((t))/I_0$ 是 $\{0\}$ 处的仿射旗簇，$\Psi:\mathcal S\to\mathrm{Perv}(I_0\backslash G((t))/I_0)=\mathrm{Perv}_{I_0}(\mathrm{Fl})$ 是 Gaitsgory 的邻近闭链函子 [16]，$V_\varphi(K):=R\Gamma_c(\mathrm{Fl}_G,K\otimes\mathrm{pr}_1^*A_\varphi)$ 如 (4.4)。为给出 ${}_\varphi\ckG$ 的钉住，需找到引理 9.3(2) 证明中 (P1) 与 (P2) 的类比。

据 Arkhipov--Bezrukavnikov [1, Theorem 4]，每个对象 $\Psi(K)$ 带 $X_*(T)$-滤链，其相伴分次片为 Wakimoto 层。更精确地，他们构造了函子
\begin{equation}
\bigoplus_{\mu\in X_*(T)}W^\mu\circ\Psi(-):\mathcal S\to\Vec^{X_*(T)}.
\end{equation}
这里对带 Wakimoto 滤链 $\Psi(K)_{\le\mu}$ 的 $\Psi(K)\in\mathrm{Perv}_{I_0}(\mathrm{Fl}_G)$，记 $\mathrm{gr}_\mu\Psi(K)=\Psi(K)_{\le\mu}/\Psi(K)_{<\mu}=J_\mu\otimes W^\mu(K)$，其中 $J_\mu$ 是 Wakimoto 层 [1, §3.2]，$W^\mu(K)$ 是带 Frobenius 作用的向量空间。在 [1, Theorem 6] 中证明了 (9.5) 是张量的。

对每个 $\mu$，设 $j_\mu:\mathrm{Fl}_\mu=I_0t^\mu I_0/I_0\hookrightarrow\mathrm{Fl}_G$ 为包含。按如下方式固定 $J_\mu$ 的 Frobenius 结构：对正则支配的 $\mu$，令 $J_\mu=j_{\mu,*}\mathbb{Q}_\ell[\langle2\rho,\mu\rangle](\langle2\rho,\mu\rangle)$；对正则反支配的 $\mu$，令 $J_\mu=j_{\mu,!}\mathbb{Q}_\ell[\langle2\rho,\mu\rangle]$；对一般的 $\mu=\mu_1+\mu_2$（$\mu_1$ 正则支配、$\mu_2$ 正则反支配），令 $J_\mu=J_{\mu_1}*^{I_0}J_{\mu_2}$（其中 $*^{I_0}$ 是 $\mathrm{Fl}_G$ 上的卷积）。由 [1, Lemma 8, Corollary 1] 知 $J_\mu$ 良定义。此归一化保证在 $J_\mu$ 的合成列中 $\delta=\mathrm{IC}_1$ 恰出现一次，重数空间为 $\mathbb{Q}_\ell$ 且作为平凡 Frobenius 模（见 [1, Lemma 3(a)]，作混合情形的显然调整）。由于除 $e_w=1$ 外所有 $\mathrm{IC}_{e_w}$ 都被 $V_\varphi$ 消去，得
\begin{equation}
V_\varphi(J_\mu)=\mathbb{Q}_\ell\ \text{作为平凡 Frobenius 模，对一切}\ \mu\in X_*(T).
\end{equation}

\medskip\noindent\textbf{断言.} 对 $K\in\mathcal S$，$W^\mu(K)$ 是权 $\langle2\rho,\mu\rangle$ 的纯模。事实上，有函子的自然同构 $h_\mu\cong W^\mu$（$h_\mu$ 是 [30, Theorem 3.5, 3.6] 中的权函子，已在引理 9.3 的证明中用到）。

\begin{proof}[断言的证明]
先回忆上引文献中权函子的定义。对每个 $\mu\in X_*(T)$，设 $S_\mu\subset\mathrm{Gr}_G$ 是含 $t^\mu$ 的 $U((t))$-轨道。上引文献中定义的权函子为 $h_\mu(K)=H^*_c(S_\mu,K)$，它集中在次数 $\langle2\rho,\mu\rangle$。

设 $\pi:\mathrm{Fl}_G\to\mathrm{Gr}_G$ 为投影，则 $\pi^{-1}(S_\mu)=\bigsqcup_{w\in W}\widetilde S_{\mu w}$，其中 $\widetilde S_{\mu w}\subset\mathrm{Fl}_G$ 是含 $t^\mu w$ 的 $U((t))$-轨道。对 $K\in\mathcal S$ 有自然同构
\[
h_\mu(K)=H^*_c(S_\mu,\pi_!\Psi(K))\cong H^*_c(\pi^{-1}(S_\mu),\Psi(K))\xrightarrow{\sim}H^*_c(\widetilde S_\mu,\Psi(K))\xrightarrow{\sim}W^\mu(K).
\]
这里第一个等式来自 $\pi_!\Psi(K)=K$，最后两个同构来自 [1, Theorem 4(2)]（须对 Frobenius 结构多加小心）。

由于 $h_i(K)=H^i(\mathrm{Gr}_G,K)$ 是权 $i$ 的纯模，而 $h_\mu(K)$ 是 $h_{\langle2\rho,\mu\rangle}$ 的直和项，故 $h_\mu(K)$ 是权 $\langle2\rho,\mu\rangle$ 的纯模。因此 $W^\mu(K)$ 也是权 $\langle2\rho,\mu\rangle$ 的纯模。
\end{proof}

下面构造自然同构 $\bigoplus_\mu W^\mu\cong\omega_\varphi$。对每个 $K\in\mathcal S$，$\Psi(K)$ 上的 Wakimoto 滤链给出计算 $V_\varphi(\Psi(K))$ 的谱序列，其 $E_1$ 页为 $V_\varphi(\mathrm{gr}_\mu\Psi(K))$。由 (9.6)，
\begin{equation}
V_\varphi(\mathrm{gr}_\mu\Psi(K))=V_\varphi(J_\mu)\otimes W^\mu(K)=W^\mu(K)
\end{equation}
集中在次数 $0$，故谱序列在 $E_1$ 处退化。谱序列的极限给出 Frobenius 模 $\omega_\varphi(K)=V_\varphi(\Psi(K))$ 上的 Wakimoto 滤链，记作 $w_{\le\mu}$。由断言与 (9.7)，该滤链细化权滤链 $w_{\le i}$：
\[
\mathrm{gr}^w_i\omega_\varphi(K)=\bigoplus_{\langle2\rho,\mu\rangle=i}\mathrm{gr}^w_\mu\omega_\varphi(K).
\]
由于 $\omega_\varphi(K)$ 是 $\mathbb{Q}_\ell$-向量空间中的 Frobenius 模，权滤链典范分裂。因此 Wakimoto 滤链 $w_{\le\mu}\omega_\varphi(K)$ 也典范分裂。这给出典范同构 $\bigoplus_\mu W^\mu\cong\omega_\varphi$。

最后，主幂幺元由单值作用（$0$ 处驯惯性群的拓扑生成元）在邻近闭链上的对数给出（[16, Theorem 2]）：$M_K:\Psi(K)\to\Psi(K)(-1)$。

$S^1_\varphi(k)$ 的作用与邻近闭链函子交换，故与单值 $M_K$ 交换。它置换 Wakimoto 层，从而经 $S^1_\varphi(k)\to\operatorname{Aut}^\dagger(G)$ 在 $X_*(T)$ 上的作用置换函子 $W^\mu$，并保持 $\langle2\rho,\mu\rangle$。因此 $S^1_\varphi(k)$-作用保持钉住 $\ddagger$。
\end{proof}
